\begin{cabstract}
随着具身智能与泛在智能物体逐步进入生产和生活环境，人机关系正由人使用单一设备完成离散任务，转向人与多个具身智能体共享物理情境、持续感知彼此状态并动态协作的人机共生。在这一过程中，交互对象、任务进程与环境状态均可能实时变化，用户因而不仅需要“操作设备”，还需要在多个共处对象之间迅速辨认当前目标、自然地表达交互意图，并将指令准确地作用于特定对象。以协同工厂为例，操作人员可能需要在机械臂、移动机器人和智能工位之间连续切换，查看某一对象的状态并随即下达指令，而不应反复退出当前任务以搜索设备。由此，面向多对象的、符合人类空间认知习惯的直观定向交互，成为人机共生时代的新需求。

然而，现有交互方式难以满足这一需求。二维码扫描、设备列表检索和无线配对等方式忽视了指向、朝向与注视中已经包含的空间关注信息；即使用户已明确关注某一对象，仍需在界面中重新查找、确认并连接设备，使物理对象选择、数字身份识别、连接建立与动作执行相互割裂。究其根因，现有系统无法将人的自然注意力直接转化为可靠的对象绑定，只能借助额外的显式步骤恢复交互过程中丢失的注意力信息，因而增加了多对象切换的时间、操作与认知负担。因此，如何捕捉并利用可测量的注意力，在不增加额外操作的前提下准确绑定交互对象，是实现直观定向交互的核心问题。
本文围绕\textbf{如何挖掘、利用人真实注意力并可靠地转化为高效、精准的定向交互}这一问题展开研究。通过对人类注意力信息由浅入深的挖掘和利用，形成\textbf{“所指即所控—所看即所控—所思即所控”}的三种定向交互模式，旨在为未来人机共生提供一个新的交互模式解决方案。具体来说，主要研究工作如下：

（1）针对广泛普及的手机等手持终端，本文提出基于指向的非接触定向交互系统 RetroAR。本文首先发现，手机指向所携带的用户位置和方向信息可由可见光的定向传播特性自然保留。基于此，设计了面向商用手机的可见光定向交互机制：用户以手机闪光灯照射逆反射标签，标签调制反射光并由手机后置摄像头接收，无须检索和连接无线设备。为兼顾数据通信与空间锚定，本文进一步设计动态光标签 ViTag，并结合消息级动态兴趣区、快速凸包检测和自适应时空混合编解码，使系统在手机运动、视角变化和复杂光学背景下稳定读取标签并估计六自由度位姿。系统评估表明，RetroAR 最远可在约 $4\,\mathrm{m}$ 距离和约 $100^\circ$ 视角下工作；用户研究表明，其交互开销相比传统Wi-Fi扫码方案降低50\%，且交互负担更小。RetroAR 验证了光学链路利用“指向”提升定向交互效率的可行性，为“所指即所控”提供系统依据；

（2）针对快速普及的AR眼镜终端，本文提出基于互补红外光学通信的定向交互系统 EchoSight。本文发现，相较于需要主动抬起并对准的手持设备，眼镜方向能够自然贴近用户的观察方向，而红外光束的定向特性可以保留用户注意力指向性，使目标识别与控制一步完成：用户看向目标并做出手势后，眼镜端 SightReader 发射调制红外光，对象端 EchoTag 接收指令并调制反射光，回送身份与状态信息，无须预先注册或建立无线连接。本文进一步设计双路互补光学前端，结合双波段差分信号、对准检测和按需唤醒机制，抑制环境光干扰与非目标响应，实现低时延双向通信。系统评估表明，EchoSight 最远可在约 $5\,\mathrm{m}$ 距离和约 $120^\circ$ 视角下工作；用户研究表明，其交互效率相比二维码扫描和语音控制方案至少提升5倍，且主观负担更低、交互流畅性更好。EchoSight 验证了“视线”比“指向”承载更直观的注意力信息，为“所看即所控”提供系统依据；

（3）针对未来小型化、多模态的新式穿戴设备，本文拟提出利用眼动捕获注意力并结合脑电解码识别意图的定向交互系统 Omnisight。本文首先观察到，用户的眼睛注视朝向和眼镜朝向往往不对齐，且注视方向真正代表了其注意力落点；基于此，设计了由眼动驱动的“物理空间三维鼠标指针”：在AR眼镜上集成基于双目相机的眼动追踪模块，并设计光学云台实现LED波束转向，使光束主动追随用户的真实视线并对准关注目标。与此同时，本文进一步引入基于水凝胶电极的脑电意图识别机制，融合眼动和任务相关脑电信息，区分“看向但不操作”与“主动选择并执行”，以减少因纯注视无交互意图所带来的误操作；Omnisight 的研究旨在探索从“方向近似注意”到“眼动对齐+意图判别”的可行路径，为“所思即所控”提供待验证的系统方案。

\end{cabstract}

\begin{eabstract}
As embodied intelligence and ubiquitous intelligent objects increasingly enter production and everyday environments, the relationship between humans and machines is shifting from the use of a single device for discrete tasks toward human--machine symbiosis, in which people and multiple embodied agents share a physical context, continuously perceive one another's states, and collaborate dynamically. In this setting, interaction targets, task progress, and environmental conditions may all change in real time. Users therefore need not only to operate devices, but also to rapidly identify the intended target among multiple co-located objects, express interaction intent naturally, and direct commands accurately to a specific object. In a collaborative factory, for example, an operator may need to switch continuously among a robotic arm, a mobile robot, and an intelligent workstation, checking the status of one object and immediately issuing a command without repeatedly leaving the current task to search for devices. Intuitive directional interaction that supports multiple objects and conforms to human spatial cognition has thus become a new requirement in the era of human--machine symbiosis.

Existing interaction methods, however, struggle to meet this requirement. QR-code scanning, device-list search, and wireless pairing overlook the spatial focus already conveyed by pointing, orientation, and gaze. Even after users have clearly attended to an object, they must still locate, confirm, and connect to the device again through an interface, separating physical-object selection, digital-identity recognition, connection establishment, and action execution. The fundamental reason is that existing systems cannot directly transform natural human attention into reliable object binding. Instead, they rely on additional explicit steps to recover attention information lost during interaction, thereby increasing the temporal, physical, and cognitive costs of switching among multiple objects. Capturing and exploiting measurable attention to bind an interaction target accurately without additional operations is therefore a central challenge in realizing intuitive directional interaction.

This dissertation investigates \textbf{how to uncover and exploit genuine human attention and reliably transform it into efficient and accurate directional interaction}. By progressively exploring and utilizing human attention, it develops three directional interaction paradigms: \textbf{``point to control, look to control, and intend to control.''} The goal is to provide a new interaction paradigm for future human--machine symbiosis. The main research contributions are as follows.

(1) For widely used handheld devices such as smartphones, this dissertation presents RetroAR, a pointing-based contactless directional interaction system. It first observes that the user's position and orientation conveyed by smartphone pointing can be naturally preserved by the directional propagation of visible light. Based on this observation, it designs a visible-light directional interaction mechanism for commodity smartphones: the user illuminates a retroreflective tag with the smartphone flashlight, while the tag modulates the reflected light and the rear camera receives it, eliminating the need to search for and connect to wireless devices. To support both data communication and spatial anchoring, this dissertation further designs the dynamic optical tag ViTag and combines a message-level dynamic region of interest, fast convex-hull detection, and adaptive spatiotemporal hybrid coding and decoding. These techniques enable the system to read tags reliably and estimate six-degree-of-freedom poses under smartphone motion, viewpoint changes, and complex optical backgrounds. System evaluations show that RetroAR operates at distances of up to approximately $4\,\mathrm{m}$ and viewing angles of up to approximately $100^\circ$. A user study shows that it reduces interaction overhead by 50\% compared with a conventional Wi-Fi and QR-code-based approach, while imposing a lower interaction burden. RetroAR demonstrates the feasibility of using an optical link to improve directional interaction efficiency through pointing, providing a system-level basis for ``point to control.''

(2) For the rapidly growing class of augmented-reality glasses, this dissertation presents EchoSight, a directional interaction system based on complementary infrared optical communication. Compared with handheld devices that must be actively raised and aimed, the orientation of glasses more naturally follows the user's viewing direction, while the directionality of an infrared beam preserves the direction of the user's attention, allowing target identification and control to be completed in a single step. After the user looks at a target and performs a gesture, the glasses-mounted SightReader emits modulated infrared light; the object-mounted EchoTag receives the command, modulates the reflected light, and returns identity and status information without preregistration or a wireless connection. This dissertation further designs a dual-path complementary optical front end and combines dual-band differential signaling, alignment detection, and on-demand wake-up to suppress ambient-light interference and responses from non-target objects, thereby enabling low-latency bidirectional communication. System evaluations show that EchoSight operates at distances of up to approximately $5\,\mathrm{m}$ and viewing angles of up to approximately $120^\circ$. A user study shows that its interaction efficiency is at least five times that of QR-code scanning and voice-control approaches, with a lower subjective workload and greater interaction fluency. EchoSight demonstrates that looking conveys attention more intuitively than pointing, providing a system-level basis for ``look to control.''

(3) For future miniaturized and multimodal wearable devices, this dissertation proposes Omnisight, a directional interaction system that captures attention through eye tracking and recognizes intent through EEG decoding. It first observes that the direction of a user's gaze is often misaligned with the orientation of the glasses and that gaze direction more faithfully indicates where the user's attention is focused. Based on this observation, it designs an eye-tracking-driven ``three-dimensional mouse pointer in physical space.'' A binocular-camera-based eye-tracking module is integrated into AR glasses, and an optical gimbal steers an LED beam so that it actively follows the user's actual line of sight and aligns with the attended target. Meanwhile, this dissertation introduces an EEG-based intent-recognition mechanism using hydrogel electrodes. By combining eye-tracking data with task-related EEG signals, the mechanism distinguishes between ``looking without acting'' and ``actively selecting and executing,'' thereby reducing unintended operations caused by gaze without interaction intent. The Omnisight study aims to explore a feasible path from approximating attention through device orientation to aligning with gaze and discriminating intent, providing a system design to be validated for ``intend to control.''
\end{eabstract}
